\section{Security of \system}
\label{sec:security}

\begin{theorem}
The \onehop\ operator in Algorithm~\ref{alg:onehop} is oblivious.
\end{theorem}

\begin{proof}
Control-flow obliviousness of the \onehop\ operator is immediate from Algorithm~\ref{alg:onehop}, as the entire algorithm has no input-dependent branching statements. 
So, all that is left is to prove the obliviousness of the memory-access pattern of \onehop. 

To this end, we provide a simulator $\mathcal{S}_{\onehop}$, that takes as input the acceptable leakage
$\mathcal{L}$ and produces as output a trace of memory accesses induced $\mathcal{S_{\onehop}}(\mathcal{L})$, which is indistinguishable from the $\mathsf{Trace}_{\mathsf{Alg}}$, the memory access pattern trace induced by a real execution of the \onehop\ algorithm.


For \onehop\, we define $\mathcal{L} = \{ |S|, |D|, |E|, |R|\}$, where
$|S|$ and $|D|$ are the sizes of the source and destination nodes list respectively,
$|E|$ is the size of the edge list, and
$|R|$ is the size of the output.
The leakages are all public parameters corresponding to the size of each database, and the final result size.
In the case where \onehop\ is used for an intermediate computation, then the $|R|$ can be ommitted in the leakage description $\mathcal{L}$ as it will be the same as $|E|$; as we do not leak intermediate sizes. 
The simulation is identical in both cases modulo the final output returned by the simulator.

Simulators for the basic oblivious primitives of \textsc{OCmpSet}, \textsc{OSort}
($\mathcal{S_{\text{OSort}}}$), \textsc{OCompact}~\cite{sasy2022fast}$(\mathcal{S_{\text{OCompact}}})$, and for an oblivious hash table ($\mathcal{S}_{\text{OHT}})$~\cite{zheng2025h2o2ram} are already well-defined in prior work.
Simulating linear passes over an array while performing any oblivious operator is also trivial, we refer to as $\mathcal{S}_{\text{LP}}$; i.e. $\mathcal{S}_{\text{LP}}$ immediately captures simulating the reduplication and deduplication steps involved in the \onehop\ algorithm.
Our simulator for \onehop\ makes black-box use of the aforementioned simulators for fundamental primitives.

The \onehop\ simulator can be constructed using these building-block simulator primitives.
Lines 2-5 (probing the source side) and 7 (the symmetric equivalent for the destination side) is simulated by simply invoking $\mathcal{S}_{\text{LP}}$, $\mathcal{S}_{\text{OHT}}$, followed by $\mathcal{S}_{\text{LP}}$ again. Note that all the above simulators do not take any real input but simply takes as input the dimensions of S, D, and E; which are all supplied to $\mathcal{S_{\onehop}}$ via $\mathcal{L}$.
Finally steps 8 and 12 are simulated by invoking $\mathcal{S_{\text{OSort}}}$ and $\mathcal{S_{\text{OCompact}}}$.
$\mathcal{S_{\onehop}}$ concatenates and returns the traces produced by the above steps.
\end{proof}